%! Author = k
%! Date = 2024/11/20

% Preamble
\documentclass[11pt]{article}

\usepackage{amsmath}
\usepackage[UTF8]{ctex}
\usepackage{multirow}
\usepackage{booktabs} % 用于更好看的表格线
\usepackage{array} % 用于调整列宽
\usepackage{caption} % 用于设置表格标题
\usepackage{geometry}

\pagestyle{empty} % 隐藏页眉页脚

\captionsetup[table]{labelformat=empty} % 禁用表格编号

\geometry{a4paper,paperwidth = 30cm, paperheight = 30cm}

\begin{document}

    % 表格代码开始
    \begin{table}[h]
        \centering
        \renewcommand{\arraystretch}{1.5} % 调整行间距
        \setlength{\tabcolsep}{12pt} % 调整列间距
        \begin{tabular}{>{\centering}p{1.5cm}|>{\centering}p{1.5cm}|>{\centering}p{1.5cm}|>{\centering}p{1.5cm}|c} % 控制边框：| 表示有竖线，c 表示居中对齐
            \hline % 表格顶部边框
            \textbf{$R$} & \textbf{$S$} & \textbf{$Q^{n}$} & \textbf{$Q^{n + 1}$} & \textbf{触发器状态说明} \\ % 表头
            \hline % 表头下边框
            0 & 0 & 0 & 0 & \multirow{2}{*}{\centering 状态不变} \\ % 第一行
            0 & 0 & 1 & 1 &  \\ % 第二行
            \hline
            1 & 0 & 0 & 0 & \multirow{2}{*}{\centering 置0} \\ % 第三行，合并说明单元格
            1 & 0 & 1 & 0 & \\
            \hline
            0 & 1 & 0 & 1 & \multirow{2}{*}{\centering 置1} \\ % 第四行
            0 & 1 & 1 & 1 & \\
            \hline
            1 & 1 & 0 & 不定 & \multirow{2}{*}{\centering $\overline{R},\,\overline{S}$端触发高电平同时撤除后状态不定} \\ % 第五行
            1 & 1 & 1 & 不定 & \\
            \hline
        \end{tabular}
        \caption{表1 图P9.1中的或非门构成的基本RS触发器特性表} % 表格标题
        \label{tab:D_flip_flop} % 表格引用标签
    \end{table}
    % 表格代码结束

\end{document}

